\section{Testing and documentation}

\subsection{Introduction to testing}

\textit{"How do I make sure my program is correct?"}. The short answer is: in real life you don't, especially if your code is asynchronous \footnote{\textit{Asynchronous programming} allows a program to initiate a long-running task (such as an I/O operation) and remain responsive to other events while waiting.}. That is not the same as saying there is nothing you can do. For compiled languages the compiler will flag all obvious and non-obvious mistakes. This doesn't really apply to Python, since Python is interpreted, although the interpreter will stop upon syntax errors.

Besides paying attention, there are two things that you can do even in interpreted languages:
\begin{enumerate}
    \item \textit{Unit testing}
    \item \textit{Static analysis}
\end{enumerate}
Generally people hate both, but they should come right next to version control in your work-flow toolbox.

\subsection{Unit testing in practice}

While programming you should break up your code in many small pieces (or units) which should encapsulate a well-defined and (possibly) simple functionality. This is usually accomplished by means of a sensible hierarchy of functions and classes is typically the hardest task when structuring your code. Also, your code will evolve with time, so you will find yourself \textbf{refactoring code} from time to time. Remember to be DRY: don't repeat yourself.

\textbf{Unit testing} consists in making sure that each single unit is correct by implementing a series of basic checks. You know what each elementary piece of code is suppose to be so make sure it does and, more specifically, make sure it does \textit{with any valid input}! This is much simpler that testing the whole program at once (although you have to do that, too).

Unit testing is so important that many programmers embrace a practice called \textit{Test-Driven Development (TDD)}:
\begin{enumerate}
    \item Write an empty placeholder for your new function;
    \item Write all the unit tests (they will fail);
    \item Implement your function and tweak it until all the tests pass:
\end{enumerate}
You don't need to do that, but imagine how fundamental unit testing is if tests can come before the code itself!

Take a look at this simple example:

\input{pygments/unit_test_naive}


This is fine, but everything happens manually. You have to run the script yourself and you have to inspect the output yourself. As your code grows in complexity, this is not very effective. Thankfully, Python has it's own module for unit testing: \texttt{unittest}. Here's how it works:

\input{pygments/unit_test}

This is much better! The base \texttt{TestCase} class offers all the goodies for unit testing (\texttt{assertTrue()}, \texttt{assertFalse()}, \texttt{assertEqual()}, \texttt{assertAlmostEqual()} etc\ldots).

The execution can be easily made automatic: put all your unit test modules into a test folder, run \texttt{python -m unittest discover} (or, even better, write a small Makefile or .bat script to do that and that's it: all your tests will run in sequence.

Did you just find a bug in your code? Make sure you add a unit test along with the fix, so that you'll never be hurt again by that particular bug. Are you adding a new feature? Make sure the new code is covered by unit tests. You should not be obsessed by the coverage, but you should definitely aim for it to be as large as possible.

When you use \texttt{git} you must remember this: \textbf{you should always make sure that all the unit tests are passing before merging stuff on the master}. More about this in a bit when we'll be talking about continuous integration.

As usual, take a look at this URL and focus on how \texttt{git} and GitHub can handle unit testing: \url{https://lucabaldini.github.io/metarep/setup.html}.

\subsection{Static code analysis}
By its very nature, Python will show you all the errors at runtime. Say you have a bug in a part of the code that is exercised very rarely, and not covered by unit tests. Python might crash the first time you exercise it or it might happily do something that is not what you intended. It might take years for even realizing that there is a bug.

Many common mistakes can be found by just looking at the code, and in fact all of them can, at least in principle. \textbf{Static code analysis} is just that: searching for errors \textit{before} running the code. Part of it can be done programmatically. Generally, a program will not \textit{understand} your program, but a program can be trained to spot some kind of errors and inconsistencies. \texttt{Pylint} and \texttt{pyflakes} are good examples of such tools. 

Let's say that you want to inspect this code using \texttt{Pylint}.

\input{pygments/linting1}

You will run \texttt{pylint file\_name.py} and you will see the following output in the terminal. As you can see, \texttt{Pylint} finds your mistakes and tells you where you've made them. It also gives a mark to your code just to humiliate you a little more.

{\scriptsize
\begin{verbatim}
[lbaldini@nbbaldini latex]$ pylint snippets/linting1.py 
************* Module snippets.linting1
snippets/linting1.py:1:0: C0111: Missing module docstring (missing-docstring)
snippets/linting1.py:1:0: C0103: Constant name "x" doesn't conform to UPPER_CASE
                          naming style (invalid-name)
snippets/linting1.py:2:0: C0103: Constant name "y" doesn't conform to UPPER_CASE
                          naming style (invalid-name)
snippets/linting1.py:3:0: C0103: Constant name "very_uncommon_condition" doesn't
                          conform to UPPER_CASE naming style (invalid-name)
snippets/linting1.py:5:14: E0602: Undefined variable 'z' (undefined-variable)

--------------------------------------------------------------------
Your code has been rated at -5.00/10 (previous run: -5.00/10, +0.00)
\end{verbatim}
}

\textit{You should consider using static code analysis routinely}. Static analysis tools tend to be quite verbose, and often times verbose is the same as annoying. They try and enforce many different good things at once, such as:
\begin{itemize}
    \item formal correctness;
    \item efficiency;
    \item avoiding anti-patterns;
    \item style guides;
    \item generic conventions.
\end{itemize}
They also are typically highly customizable. i.e., you can mute errors you don't care about. But be advised: most of the times you should probably care. Finding a good balance is generally not too hard and it will surely help you in the long run.

One more thing: recent Python 3 versions support type annotations. \textit{The Python interpreter recognizes but does nothing with annotations}, but they can be handy as comments are. The code is easier to read and even more checks with respect to an un-annotated code can be done by tools such as \texttt{mypy}.

\input{pygments/type_annotation}

\subsection{Continuous integration}
Wouldn't it be nice if somebody run all the unit tests of my package every time I push on the master or make a pull request? And, since we are at it, sent me an email if any of the tests fail? Well, such a thing exists and it is standard practice in code development. People even made a name for it: \textbf{Continuous Integration (CI)}.

CI cloud-base services exists just like code-hosting services exist. Travis-CI and circleci are two good examples. They interoperates seamlessly with GitHub, gitlab or bitbucket. Setting up CI for your package is usually fairly simple, you can find out how to do it (you guessed it) here: \url{https://lucabaldini.github.io/metarep/setup.html}. \\ \textbf{One-sentence summary: always implement continuous integration!}

\subsection{Documentation and Sphynx}

\begin{figure}[H]
\centering
\includegraphics[width=0.6\textwidth]{curve_fit_docs}
\end{figure}

Take a look at the page containing the documentation of the \texttt{curve\_fit} function from the \texttt{scipy.optimize} module. How can it be automatically generated? If you open the source code you'll notice that the documentation is written directly as a docstring in Python.

\begin{figure}[H]
\centering\includegraphics[width=0.6\textwidth]{curve_fit_code}
\end{figure}

\textbf{Sphinx} is the documentation tool for Python. You must learn to use it and love it. 

\begin{figure}[H]
\centering
\includegraphics[width=0.6\textwidth]{sphinx}
\end{figure}

Sphynx automatically processes all the relevant information to produce several types of output, most notably html and LaTeX. It considers two different sources of information:
\begin{enumerate}
    \item The doctrings in the Python modules;
    \item The additional markup files (written in \texttt{reStructuredText}) containing auxiliary information.
\end{enumerate}
Use \texttt{sphinx-quickstart} once when you setup your project and then weak the generated \texttt{conf.py} file to suit your needs. You probably don't want to see this link anymore, but here you can learn how to use Sphinx in practice: \url{https://lucabaldini.github.io/metarep/setup.html}.

Sphinx is \textit{very} powerful: the Python documentation itself (\url{https://docs.python-guide.org/}) is written in Sphinx, and so is any module documentation.

\textit{"Ok, I have the documentation compiled. Now what do I do with it?"}
Wouldn't it be nice if the documentation was automatically compiled and uploaded on the web each time you push on the master branch? This is possible and is called readthedocs.com. And, again, this is a cloud-based service that can interoperate easily with GitHub, gitlab or bitbucket. How do you do it in real life? Well, here's the link again: \url{https://lucabaldini.github.io/metarep/setup.html}.

\subsection{References}
\begin{itemize}
    \item \url{https://docs.python.org/3/library/unittest.html}
    \item \url{https://www.pylint.org/}
    \item \url{https://pypi.org/project/pyflakes/}
    \item \url{http://mypy-lang.org/}
    \item \url{https://circleci.com/}
    \item \url{https://travis-ci.org/}
    \item \url{http://www.sphinx-doc.org/en/master/}
\end{itemize}